%%%%%%%%%%  scan idx 343  |  folio 330  %%%%%%%%%%
\textbf{Définition 3.9.} Soit \(K\) un \(\lambda\)-anneau~; on appelle \(K\)-\(\lambda\)-algèbre augmentée la donnée d'une \(K\)-\(\lambda\)-algèbre \(A\) et d'un \(\lambda\)-homomorphisme de \(K\)-algèbres \(A\to K\)~; un homomorphisme de \(K\)-\(\lambda\)-algèbres augmentées est un \(\lambda\)-homomorphisme compatible aux augmentations.

\textbf{Exemple 3.9.1.} : Soit \((X,A)\) un topos commutativement localement annelé, et \(K^\bullet(X)\) le groupe de Grothendieck de la catégorie des faisceaux de \(A\)-modules localement libres de type fini sur \(X\). Nous avons vu que \(K^\bullet(X)\) peut être muni d'une structure de pré-\(\lambda\)-anneau (2.2.b, car tout facteur direct d'un libre de type fini est alors localement libre de type fini) et que \(H^0(X,\underline{\mathbb Z})\) est canoniquement muni d'une structure d'anneau binômial. On définit un \(\lambda\)-homomorphisme de \(H^0(X,\underline{\mathbb Z})\) dans \(K^\bullet(X)\) de la façon suivante : soit \(a\) une section du faisceau constant \(\underline{\mathbb Z}\) sur \(X\)~; elle définit une partition de l'objet final de \(X\) en une somme d'ouverts disjoints \(U_n\), la section \(a\) ayant pour valeur \(n\) sur l'ouvert \(U_n\)~; soit alors \(F\) le faisceau égal à \(A^n\) sur \(U_n\) pour \(n\ge0\), et à \(0\) sur \(U_n\) pour \(n\le0\), et soit \(\underline{F}'\) le faisceau égal à \(A^{-n}\) sur \(U_n\) pour \(n\le0\), et à \(0\) sur \(U_n\) pour \(n\ge0\)~; on associe à \(a\) l'élément \(F-F'\) de \(K^\bullet(X)\).

Il est immédiat de vérifier qu'on obtient bien ainsi un \(\lambda\)-homomorphisme de \(H^0(X,\underline{\mathbb Z})\) dans \(K^\bullet(X)\)~; on peut alors définir l'augmentation de \(K^\bullet(X)\) dans \(H^0(X,\underline{\mathbb Z})\), en associant à tout faisceau localement libre de type fini son rang, qui est une section de \(\underline{\mathbb Z}\)~; l'additivité du rang entraine qu'on définit ainsi une fonction sur \(K^\bullet(X)\), à valeurs dans \(H^0(X,\underline{\mathbb Z})\). On a donc muni \(K^\bullet(X)\) d'une structure de pré-\(\lambda\)-algèbre augmentée. C'est toujours de celle-ci qu'il sera question dans la suite.

